\documentclass[10pt,a4paper]{amsart}
\usepackage[margin=19mm]{geometry}
\usepackage{amsmath,amssymb,mathtools}
\usepackage[scr=rsfs,cal=euler]{mathalfa}
\usepackage{xfrac}
\usepackage[unicode,hidelinks,bookmarks=false]{hyperref}
\newcommand{\ie}{i.e.\ }
\newcommand{\cf}{cf.\ }
\newcommand{\resp}{respectively\ }
\newcommand{\Subsection}{\subsection}
\providecommand{\nobreakdash}{\nobreak\hbox{-}}
\setlength{\emergencystretch}{2em}
\multlinegap0pt
\usepackage{bm}
\usepackage{bbm}
\usepackage{dsfont}
\usepackage{xspace}
\usepackage{cancel}
\usepackage{mathtools}
\usepackage{amsthm}
\makeatletter
\@ifundefined{lemma}{\newtheorem{lemma}{Lemma}}{}
\@ifundefined{theorem}{\newtheorem{theorem}{Theorem}}{}
\makeatother
\DeclareMathOperator{\ad}{ad}
\title[A family of maximal subalgebras of the Lie algebra~$W\_n(K)$]{A family of maximal subalgebras of the Lie algebra~$W\_n(K)$}
\author{Y.Chapovskyi, A.Petravchuk}
\date{Source: arXiv:2602.19601v1 (CC BY 4.0)}
\begin{document}
\maketitle

\noindent\emph{Source and licence.} A family of maximal subalgebras of the Lie algebra~$W_n(K)$. Version arXiv:2602.19601v1 (CC BY 4.0), distributed under the Creative Commons Attribution 4.0 International licence, as recorded in the arXiv licence metadata for this exact version and shown on the abstract page https://arxiv.org/abs/2602.19601v1. Full rights notice: licenses/arxiv-2602.19601-CC-BY-4.0.txt.\par\smallskip

\noindent\textit{Editorial scope.} Counted unit: the Theorem whose source label is \texttt{th\_1} in the article (source0.tex lines 160--189, source label \texttt{th\_1}), delivered together with the article's own complete original proof of it (source0.tex lines 191--250, one \texttt{proof} environment of 60 lines). No mathematics is written, repaired or re-derived: statement and proof are the article's own text line for line. Required context retained uncounted: lm:grading (lines 140--142). Every definition, display and label of those lines is retained unchanged. Omitted from this package and not redistributed: 1--139, 143--159, 190--190, 251--452. Nothing inside a retained excerpt is omitted or altered: every retained body is delivered exactly as the article's lines are written, so no omission receipt of any kind accompanies this package. Citations: the retained excerpts carry no citation command at all, so no bibliography entry of the article is transcribed and no reference is redistributed. Reference adaptation: none was needed and none was made; every reference inside the delivered unit resolves inside it, so no unresolved reference or citation is produced. Printed slips kept literal: no printed word, symbol, hypothesis or number of the counted statement or of its proof was altered. Layout and class: the article's own class is \texttt{amsart} with packages including graphics, epsfig, graphicx, amssymb, color; the wrapper is the lane's pinned \texttt{amsart} editorial wrapper at 10pt on A4, which restates the theorem-like environments the retained text needs and adds the article's own macro definitions that the retained text needs (\texttt{\textbackslash ad}), with extra wrapper packages (bm, bbm, dsfont, xspace, cancel, mathtools). The delivered numbering is the unit's own. Assets: no retained span contains a figure environment, a graphics-inclusion command, a PDF-page inclusion, a table or a TikZ picture; no image file is copied into this package, the package declares no asset dependency and no image file is redistributed with it. Licence: version arXiv:2602.19601v1 of this article is distributed under the Creative Commons Attribution 4.0 International licence, as recorded in the arXiv licence metadata for this exact version and shown on the abstract page https://arxiv.org/abs/2602.19601v1. A full rights notice is delivered at licenses/arxiv-2602.19601-CC-BY-4.0.txt. Characters: every retained excerpt is pure ASCII, so the delivered engine, XeLaTeX with its default Latin Modern text font, reports no missing character.
\begin{lemma} \label{lm:grading}
	Let $L$ be a subalgebra of the Lie algebra $W_n(K)$. If $E_n \in L$, then $L$ is a graded subalgebra with respect to the standard grading on $W_n(K)$.
\end{lemma}

	\begin{theorem}\label{th_1}
	\quad
\begin{enumerate}
	\item
	Let $s$ be an integer with $1\le s\le n-1$.
	Then the set
	\[
	m_s(K)=\left\{
	\sum_{i=1}^s f_i\frac{\partial}{\partial x_i}
	+\sum_{j=s+1}^n g_j\frac{\partial}{\partial x_j}
	\mid f_i\in P_s,\ g_j\in P_n
	\right\}
	\]
	is a maximal subalgebra of the Lie algebra $W_n$.
	
	\item
	The set
	\[
	I_s=\left\{\sum_{j=s+1}^n g_j\frac{\partial}{\partial x_j}\mid g_j\in P_n\right\}
	\]
	is a polynomial ideal of the Lie algebra $m_s(K)$ and
	$I_s\simeq P_s\otimes \mathrm{Der}(K[x_{s+1},\ldots,x_n])$.
	Besides, the quotient algebra $m_s(K)/I_s$ is isomorphic to the Lie algebra $W_s$.
	
	\item $I_s$ is a unique non-trivial ideal of the Lie algebra  $m_s(K)$.
	
	\item
	The subalgebras $m_s(K)$ and $m_t(K)$ are non-isomorphic for $s\ne t$,
	$s,t=1,\ldots,n-1$.
\end{enumerate}\end{theorem}

\begin{proof}[Proof of (1)]
	One can immediately check that $m_s(K)$ is a subalgebra of $W_n(K)$. Let $S$ be a   subalgebra of $W_n(K)$ strictly containing the subalgebra  $m_s(K)$.
	The subalgebra $m_s(K)$ obviously contains the Euler derivation $E_n = \sum_{i=1}^n x_i \frac{\partial}{\partial x_i}$, so 
	by Lemma \ref{lm:grading} the subalgebra $S$ is graded, $S = \bigoplus_{i \ge -1} S_i$, where  $S_i=S\cap W_n^{[i]} $ and $W_n^{[i]}$ denotes the standard grading component of the Lie algebra $W_n(K)$.
	
	Repeating some parts of the proof of Lemma \ref{lm:grading} one  can show that the set $S \setminus m_s(K)$ contains nonzero  homogeneous derivations. Fix such a derivation and denote it by $D$.
	Write $D = \sum_{i=1}^n h_i \frac{\partial}{\partial x_i}, $ where  $h_i\in P_n, i=1, \ldots , n.$ Since $D \notin m_s(K)$, there exists an index $j_0 \in \{1, \dots, s\}$ such that the coefficient $h_{j_0}$ depends on a variable $x_{i_0}, i_0 \ge s + 1,$ that is, $\frac{\partial h_{j_0}}{\partial x_{i_0}}\not =0$. 	
	Write  the polynomial $h_{j_0}$ as a polynomial in  $x_{i_0}$:
	\[
	h_{j_0} = u_0 + u_1 x_{i_0} + \dots + u_d x_{i_0}^d
	\]
	for some $d \ge 1$, where $u_t$ does not depend on $x_{i_0}$ for $t=0,\dots d$  and $u_d \neq 0.$
	
	Consider the  following homogeneous derivation:
	\[
	D_1 := 
	\left[ \dots \left[ D, \frac{\partial}{\partial x_{i_0}} \right], \dots, \frac{\partial}{\partial x_{i_0}} \right] = \sum_{l=1}^n q_l \frac{\partial}{\partial x_l}, 
	\]
	where $\partial/\partial x_{i_0}$  is applied  $d-1$ times).
	Note that the coefficient $ q_{j_0} $  can be written in the form $q_{j_0} = v_0 + v_1 x_{i_0}$ with $v_1 \neq 0$, and the polynomials $v_0, v_1$ do not depend on $x_{i_0}$.
	Clearly, $D_1 \in S \setminus m_s(K)$ as well.
	Without loss of generality, we may assume that $v_1 \in K^*$. Indeed, if $v_1 \notin K^*$,  consider
	a  monomial $x_1^{l_1} \dots x_{i_0}^0 \dots x_n^{l_n}$ of maximal total degree  in the polynomial $v_1$. Acting by the linear operator 
	$$\ad(\frac{\partial}{\partial x_1})^{l_1} \dots \ \ad(\frac{\partial}{\partial x_n})^{l_n}, \  l_i\geq 0, \  i=1, \ldots , n$$ on $D_1 $
	we can obtain a derivation satisfying the condition $v_1 \in K^*$.
	Since $D_1$ is  homogeneous, it must be a linear homogeneous derivation. Acting by the operator $\ad(x_{i_0} \frac{\partial}{\partial x_{i_0}})$
	on the derivation $D_1$ we obtain the derivation 
	\[
	D_2:= \left[x_{i_0} \frac{\partial}{\partial x_{i_0}}, D_1 \right] = \sum_{i=1}^n c_i x_{i_0} \frac{\partial}{\partial x_{i}}
	=  x_{i_0}\sum_{i=1}^n c_i \frac{\partial}{\partial x_{i}},
	\]
	where $c_i \in K$, $c_{j_0} \ne 0$.
	Again, $D_2 \in S \setminus m_s(K)$. We may assume that $c_i = 0$ for all $i > s$, since 
	$\sum_{i=s+1}^n c_i x_{i_0} \frac{\partial}{\partial x_{i}} \in m_s(K)$.
	Note that
	\[
	\left[x_{j_0} \frac{\partial}{\partial x_{i_0}}, D_2 \right] = 
	\left[ x_{j_0} \frac{\partial}{\partial x_{i_0}}, x_{i_0}\sum_{i=1}^{s} c_i \frac{\partial}{\partial x_i} \right] = 
	x_{j_0} \sum_{i=1}^{s} c_i \frac{\partial}{\partial x_i} - x_{i_0} c_{j_0} \frac{\partial}{\partial x_{j_0}}.
	\]
	Since the first term in the latter sum belongs to  the subalgebra $m_s(K)$ and $c_{j_0} \ne 0,$ we  conclude that 
	$x_{i_0} \frac{\partial}{\partial x_{j_0}} \in S$.
	
	Furthermore, for any $p, 1 \le p \le s $,  we have  the equality
	$$\left[ x_{j_0} \frac{\partial}{\partial x_p}, 
	x_{i_0} \frac{\partial}{\partial x_{j_0}} \right] = 
	-x_{i_0} \frac{\partial}{\partial x_p}.$$ 
	 
	So, we can find in the subalgebra $S$ all  derivations of the form $x_{i_0} \frac{\partial}{\partial x_p}$, $1 \le p \le s$.
	
	Take an arbitrary  element of the form $f(x_1, \dots, x_n) \frac{\partial}{\partial x_{i_0}}$ and let $j \le s$.
	Then 
	$$\left[ f(x_1, \dots, x_n) \frac{\partial}{\partial x_{i_0}}, 
	x_{i_0} \frac{\partial}{\partial x_{j}} \right] = 
	f(x_1, \dots, x_n) \frac{\partial}{\partial x_{j}} - 
	x_{i_0} \frac{\partial f}{\partial x_j} \frac{\partial}{\partial x_{i_0}}.$$
	Since the second term lies in $m_s(K)$, 
	we conclude that $f \frac{\partial}{\partial x_{j}} \in S$. Therefore, $S = W_n(K)$, which shows that  
	$m_s(K)$ is a maximal subalgebra of $W_n(K)$.
\end{proof}

\end{document}
